\bookchapter{Kapitel 13: Empfehlungen für Werkzeug-Apps}

\chapterintro{In diesem Kapitel stelle ich dir ausgewählte Tools für Produktivität, Gesundheit und unterwegs vor und ergänze praktische System- und Gerätefunktionen: Aufräumen und Schutz vor Belästigung mit Apps wie Qingzhuo, die eingebauten Blockier-Schalter des Handys, Wetter- und Erdbebenwarnungen; PC-Fernsteuerung mit Sunflower/ToDesk, KI-Arbeit mit Trae und anderen, Rengong Desktop, Schnellübertragung im lokalen Netzwerk mit LocalSend und Co., die geräteübergreifende Zwischenablage, Einmalpasswörter mit Aegis/Microsoft Authenticator sowie Terminplanung mit TickTick und das Scannen von Papierdokumenten.}

\booksection{Produktivität: Zeit und Aufmerksamkeit steuern}

Produktivitäts-Apps sind das „Gerüst“ des digitalen Alltags – sie helfen, verstreide Ideen, Aufgaben und Unterlagen zu ordnen. Hier Empfehlungen nach Anwendungsfall.

\booksubsection{To-do-Listen (Aufgabenverwaltung)}

\begin{itemize}
  \item TickTick Zielgruppe: Alle, die Arbeit, Privatleben, Lernen und Termine an einem Ort führen wollen; eine der Top-Empfehlungen dieses Abschnitts. Stärken: Sync über Handy/PC/Tablet/Web; Listen + Kalender + Pomodoro in einem; Natürliche Spracheingabe wie „Morgen 15 Uhr Meeting“ oder „Jeden Montag Wochenbericht“ erkennt Zeiten und Wiederholungen automatisch; Unteraufgaben, Tags, Prioritäten, Anhänge und geteilte Listen. Ideal für „Terminplanung + Aufgaben + Tagesrückblick“.
  \item TickTick noch besser nutzen: Inbox: Alles erstmal in „Heute/Inbox“ sammeln, abends auf Arbeit, Familie, Lernen verteilen.
  \item Kalenderansicht: Meetings und Deadlines in den Kalender ziehen – Konflikte sieht man sofort, klarer als in einer reinen Listenansicht.
  \item Wiederkehrende Aufgaben: Medikamente, Rechnungen, Sport als Regel anlegen – nicht jedes Mal neu eintippen.
  \item Pomodoro: Fokusmodus direkt an der Aufgabe starten und auswerten, wie viel diese Woche „wirklich erledigt“ wurde.
  \item Mit dem Handy verknüpfen: Sprachassistent mit „Erinnere mich…“ für Systemerinnerungen; wichtige Termine zusätzlich im Systemkalender absichern.
\end{itemize}

Todoist Zielgruppe: Minimalisten, plattformübergreifende Nutzer. Stärken: Aufgeräumte Oberfläche, klare Projektstruktur (Aufgabe Unteraufgabe Unter-Unteraufgabe), Karma-Punkte als Anreiz für Gewohnheiten.

Microsoft To Do Zielgruppe: Büro- und Outlook-Nutzer. Stärken: Komplett kostenlos, werbefrei. „Mein Tag“ wird täglich geleert und entlastet den Kopf. Tiefe Integration in das Microsoft-Ökosystem.

System-Notizen/Erinnerungen Zielgruppe: Wenignutzer ohne Drittanbieter-Apps. Stärken: Native Erfahrung, keine Verzögerung, per Siri/Xiaoai hinzufügen, Widget-Unterstützung.

\booksubsection{Pomodoro (Fokus-Training)}

\begin{itemize}
  \item Forest Zielgruppe: Wer zu oft am Handy klebt. Stärken: Während des Fokus wächst ein Baum; verlässt man die App zu früh, stirbt er ab. „Virtuelles Leben“ als positives Feedback, gemeinsames Pflanzen mit Freunden möglich.
  \item Tide – Für wen geeignet: alle, die zum Konzentrieren oder Einschlafen White Noise brauchen. Das Besondere: Naturklänge (Regen, Meereswellen, Wald) plus Pomodoro-Timer plus Meditation – alles in einer App, mit extrem minimalistischer, wohltuender Oberfläche.
  \item Prinzip: Klassisch 25 Minuten Fokus + 5 Minuten Pause, je nach Zustand auch 50/10 oder 90/20.
\end{itemize}

\booksubsection{Kalender (Zeit sichtbar machen)}

\begin{itemize}
  \item Systemkalender Zielgruppe: Alle. Stärken: Native Push-Nachrichten am stabilsten, sehr stromsparend, per Siri/Xiaoai eintragbar.
  \item Google Calendar Zielgruppe: Mehrere Geräte, gemeinsame Termine. Stärken: Die Funktion „Ziele“ plant Trainings-/Lernzeiten automatisch; Gmail erkennt Termine in Mails.
\end{itemize}

\booksubsection{Notizen (Wissen festhalten)}

\begin{itemize}
  \item Notion Zielgruppe: Wer Notizen, Aufgaben, Datenbanken und Wiki an einem Ort will; Studierende, Freelancer, Teams. Stärken: All-in-one-Wissensdatenbank, Block-Editor, persönliche Startseite, Leselisten, Projektboards, reiche Vorlagen.
  \item Obsidian Zielgruppe: Tiefes Wissensmanagement, Forschung, Schreiben. Stärken: Lokale Markdown-Dateien, bidirektionale Verlinkung + Graph – Ihr „zweites Gehirn“, die Daten gehören ganz Ihnen.
  \item flomo Zielgruppe: Kurze Einfälle festhalten, ohne lange Texte. Stärken: Kartenbasierte Notizen wie ein kurzer Post, \#Tags zum Sortieren, Zeitstrahl zum Nachschlagen.
  \item Apple Notizen Zielgruppe: iPhone/Mac-Nutzer, leichte Notizen. Stärken: Startet blitzschnell, unterstützt Dokumentenscan/Zeichnung/Checklisten/iCloud-Sync, kostenlos und tief integriert.
  \item Evernote Zielgruppe: Vielsammler von Web-Clips und Archiven. Stärken: Starkes Clipper-Plugin, OCR für Bildtexte, bewährt und stabil – Premium jedoch kostenpflichtig.
\end{itemize}

\booksubsection{PDF-Werkzeuge}

\begin{itemize}
  \item PDF Expert Zielgruppe: iPad/iPhone-Nutzer, Business-Anwender, Studierende. Stärken: Gilt als stärkster PDF-Editor auf iOS, flüssige Annotation, Signatur, Zusammenfügen/Trennen, Formulare ausfüllen.
  \item WPS Zielgruppe: Office-Nutzer, die Word/Excel/PPT/PDF parallel brauchen. Stärken: Alles in einer Suite, gute Kompatibilität mit chinesischen Cloud-Laufwerken und WeChat-Dateien; die Gratisversion reicht für den Alltag.
  \item Adobe Acrobat Zielgruppe: Hohe Anforderungen an PDF-Standards (Recht, Design, Publishing). Stärken: Offizieller Standard, zuverlässige PDF/A-Konvertierung, geräteübergreifende Sync und professionelle E-Signatur.
\end{itemize}

\booksubsection{Scannen (Papier digitalisieren)}

Papierverträge, Rechnungen, Hausaufgaben, Whiteboard-Notizen, Ausweiskopien … – die Handykamera ersetzt inzwischen die meisten Scanner. Ablauf: Sauber fotografieren Ränder automatisch korrigieren PDF/Bild exportieren optional OCR.

\begin{itemize}
  \item CamScanner Zielgruppe: Geschäftsleute und Studierende, die oft Verträge, Belege oder Ausweise scannen. Stärken: Automatische Kanten, Perspektivenkorrektur, OCR, mehrseitige PDFs mit Cloud-Sync.
  \item Microsoft Lens Zielgruppe: Studierende, Whiteboard-/Buchseiten in Word umwandeln. Stärken: Kostenlos, werbefrei, Export nach OneNote/Word per Tap, hohe Texterkennungsqualität.
  \item Apple Notizen Scan Zielgruppe: iPhone-Nutzer, leichte Scan-Bedürfnisse. Stärken: Keine extra App nötig – in Notizen auf das Kamerasymbol tippen, automatische Schattenentfernung und Signatur, kostenlos und praktisch.
  \item Baimiao / CamScanner Web / „Dokumentenscan“ der Kamera-App: Huawei-, Xiaomi-, OPPO- und vivo-Kameras haben meist einen „Dokument/Scan“-Modus. Unter Android auch Google Camera oder der System-Scanner. Gut, um schnell eine Vertragsseite zu knipsen.
  \item Scan-Tipps: Gleichmäßiges Licht, Handy parallel zum Papier, auf die vier Ecken achten. Mehrseitige Dokumente in einer Reihenfolge fotografieren und dann als PDF zusammenfügen. Sensible Ausweise mit Wassermarke und sicherem Speicherort behandeln, danach endgültig aus „Zuletzt gelöscht“ entfernen.
\end{itemize}

\booksection{Gesundheit \& Sport}

Gesundheits-Apps sind die unterschätzten „Privatassistenten“ des Smartphones – sie tracken Schlaf, Herzfrequenz und Bewegung dauerhaft und zeigen den Langfrist-Zustand oft besser als ein seltener Arztbesuch.

\booksubsection{Fitness}

\begin{itemize}
  \item Keep Zielgruppe: Home-Training, Anfänger, Fetaufbau. Stärken: Sehr viele kostenlose Kurse von HIIT bis Yoga, KI-generierte Trainingspläne, lebendige Community.
  \item Nike Training Club (NTC) Zielgruppe: Wer professionelle Trainerleitung und systematisches Training sucht. Stärken: Offiziell von Nike, alle Kurse kostenlos, Übungen von Profis vorgeführt, klar nach Muskelgruppe/Gerät sortiert.
  \item Apple Fitness Zielgruppe: Apple-Watch-Nutzer, die über Ringe motiviert werden wollen. Stärken: Drei tägliche Aktivitätsringe (Bewegung/Training/Stehen), visuelle Ziele, nahtlose Herzfrequenzdaten der Uhr.
\end{itemize}

\booksubsection{Schlaf-Tracking}

\begin{itemize}
  \item AutoSleep Zielgruppe: Apple-Watch-Nutzer, die Tief- und Leichtschlafanteile kennen wollen. Stärken: Automatisches Tracking im Schlaf, kein manuelles Starten/Stoppen, Schlafqualitäts-Score und Einschlaf-Tipps.
  \item Sleep Cycle Zielgruppe: Ohne Uhr schlafen und mit dem Handy messen. Stärken: Analysiert über das Mikrofon Bewegungen und Schnarchen, weckt in der Leichtschlafphase – aufstehen fällt leichter.
  \item Xiaomi Sports \& Health Zielgruppe: Mi-Band-/Watch-Nutzer. Stärken: Kostenlos, tief integriert mit Mi Band, Herzfrequenz/Blutsauerstoff/Schlaf in einem, gutes Preis-Leistungs-Verhältnis.
\end{itemize}

\booksubsection{Meditation \& Entspannung}

\begin{itemize}
  \item Calm Zielgruppe: Gestresste Berufstätige, Einschlafprobleme. Stärken: „Daily Calm“-Geführte Meditation + Sleep Stories (von Prominenten gesprochen), schöne Oberfläche, reichhaltige englische Inhalte.
  \item Tide – Für wen geeignet: chinesischsprachige Nutzer, die zum Konzentrieren oder Einschlafen White Noise brauchen. Das Besondere: ein Qualitätsprodukt aus China mit Naturklängen, Meditationskursen und Pomodoro-Timer – die chinesischen Anleitungen fühlen sich besonders vertraut an.
  \item Little Sleep Zielgruppe: Einschlafhilfen, Schnarch-/Traumreden aufzeichnen. Stärken: Sehr viele Geräusche kombinierbar (Regen + Ventilator + Lagerfeuer), Sleep-Timer, Traumreden-Aufnahmen sind praktisch und amüsant.
\end{itemize}

\booksubsection{Frauengesundheit}

\begin{itemize}
  \item Flo Zielgruppe: Zyklus, Kinderwunsch, Frauengesundheit. Stärken: KI-Prognose für Periode/Eisprung recht genau, anonyme Community-Fragen, ansprechendes Design.
  \item Meiyou – Für wen geeignet: Frauen in China und alle, die chinesischsprachige Communities und Aufklärungscontent suchen. Das Besondere: reichlich lokalisierte Inhalte, mit passendem Wissen für jede Phase – Zyklus, Kinderwunsch, Schwangerschaft und Erziehung.
\end{itemize}

\booksubsection{Medikamenten-Erinnerung}

\begin{itemize}
  \item Medication Assistant Zielgruppe: Chronisch Kranke, Angehörige, die ältere Menschen bei der Medikation unterstützen. Stärken: Medikamente per Scan erfassen, Liste automatisch anlegen, Erinnerungen, Einnahme-Dokumentation zum Teilen mit der Familie.
\end{itemize}

\booksubsection{Datenzentrale}

\begin{itemize}
  \item Apple Health / Google Fit Zielgruppe: Alle Smartphone-Nutzer. Stärken: Bündelt Daten aus Handy, Band und Drittanbieter-Apps (Schritte, Herzfrequenz, Schlaf, Zyklus usw.) in einer „Gesundheitsakte“.
\end{itemize}

\booksection{Reise \& Mobilität}

Unterwegs sparen gute Apps viele Stolpersteine. Hier Empfehlungen entlang der gesamten Reise.

\booksubsection{Flug \& Bahn}

\begin{itemize}
  \item Umetrip – Für wen geeignet: Vielflieger im Geschäftsleben. Das Besondere: offizielle Luftfahrt-Daten, bei Flugstatus, Gate-Änderungen und Verspätungswarnungen die verlässlichste Quelle; auch Flugzeugtyp und Sitzplan sind abrufbar.
  \item 12306 Zielgruppe: Alle Bahnreisenden in China. Stärken: Einzige offizielle Ticketplattform, Wartelisten-Kauf (bei Stornierungen automatisch zuschlagen), Studenten-/Gruppentickets bequem.
  \item VariFlight – Für wen geeignet: alle, die auf Pünktlichkeit achten und jemanden zum Abholen brauchen. Das Besondere: treffsichere Verspätungsprognosen; du kannst deine Reise mit der Familie teilen, die deinen Standort und deine Ankunftszeit live sehen kann.
\end{itemize}

\booksubsection{Karte \& Navigation}

\begin{itemize}
  \item Amap (Gaode Maps) – Für wen geeignet: Nutzer in chinesischen Städten, Autofahrer und alle, die per App ein Taxi rufen. Das Besondere: genaue Verkehrslage in Echtzeit und starke Routen, die Staus vermeiden; Taxi und Sammelfahrt an einem Ort, dazu viele Umsteige-Optionen im ÖPNV.
  \item Baidu Maps Zielgruppe: Ortssuche, Indoor-Navigation. Stärken: Street View / AR-Navigation, detaillierte Indoor-Karten für Malls und Sehenswürdigkeiten, viele Sprachpakete.
  \item Google Maps Zielgruppe: Auslandsreisen. Stärken: Weltweit beste Abdeckung, authentische lokale Bewertungen, Street View und Offline-Karten sind im Ausland Gold wert.
\end{itemize}

\booksubsection{Übersetzung}

\begin{itemize}
  \item Google Translate Zielgruppe: Auslandsreisende, sprachübergreifendes Lernen. Stärken: „Konversation“ für Echtzeit-Übersetzung in beide Richtungen, „Kamera“ für Sofort-Übersetzung – ein Reise-Muss. Offline-Sprachpakete verfügbar.
\end{itemize}

\booksubsection{Währungsrechner}

\begin{itemize}
  \item Minimalistischer Währungsrechner – Für wen geeignet: Reisende, Online-Shopper im Ausland und alle, die grenzüberschreitend einkaufen. Das Besondere: extrem aufgeräumte Oberfläche, über 160 Währungen, tagesaktuelle Kurse, gleichzeitiger Vergleich mehrerer Währungen und schnelles Umrechnen durch langes Drücken auf eine Zahl.
\end{itemize}

\booksubsection{Wetter}

\begin{itemize}
  \item System-Wetter / China Weather Network Zielgruppe: Alle, erste Wahl für Prognosen und offizielle Unwetterwarnungen. Stärken: Hersteller-Apps und „China Weather“ greifen auf die Warnmeldungen der Wetterdienste zu. Starkregen, Taifun und Hitze erscheinen als Banner/Push – hohe Autorität.
  \item Caiyun Weather – Für wen geeignet: alle, die eine Regenvorhersage auf die Minute genau brauchen. Das Besondere: bei der Frage „Regnet es in den nächsten zwei Stunden?“ branchenführend genau; die Regenfront ist auf einer animierten Karte direkt sichtbar, und vor dem Regen bekommst du von selbst eine Erinnerung.
  \item Carrot Weather Zielgruppe: Wetter-Fans mit Hang zu Humor und Widgets. Stärken: Scharfzüngige KI-Stimme, viele Widgets für den Homescreen, professionelle Radaransicht; die Premium-Version kann mehr.
  \item Warnhinweise: Wetter-Apps brauchen Benachrichtigungsrechte, sonst kommen Warnungen nicht an. Es gibt die Stufen blau, gelb, orange und rot – rot bedeutet höchstes Unwetterrisiko und Reisen sollten angepasst werden. Mehr dazu im Abschnitt „Wetterwarnungen / Naturkatastrophen-Warnungen / Erdbebenwarnungen“ dieses Kapitels.
\end{itemize}

\booksubsection{Reiseführer \& Tipps}

\begin{itemize}
  \item Mafengwo Zielgruppe: Ausführliche Planung, Roadtrips, Tiefenreisen. Stärken: Viele originale User-Reiseberichte, detaillierte Aktivitäten/Routen/Fallen-Tipps, Offline-Reiseführer verfügbar.
  \item Xiaohongshu Zielgruppe: Junge Reisende, die echte Erfahrungen und Insider-Tipps suchen. Stärken: Anschauliche Bild-/Kurzvideo-Beiträge, Inhalte wie „3 Tage in XY“ sind aktuell und zeitnah, gute Suche.
\end{itemize}

\booksection{Automatisierung – erste Schritte}

Automatisierung ist ein Hebel: einmal einrichten, dauerhaft Zeit sparen – und die täglich dutzenden Male wiederholten Handgriffe schrumpfen gegen null. Dieser Abschnitt ist nur eine Begriffseinführung; konkrete Anleitungen und Tutorials finden Sie in Kapitel 23: Denken wie ein Programmierer – Einführung in die Automatisierung.

\begin{itemize}
  \item iOS „Kurzbefehle“: Apples offizielles Automatisierungswerkzeug, ganz ohne Programmieren. Mit einem Tap lassen sich Kombinationen auslösen wie „Zuhause: WLAN an + Helligkeit runter + Wetter vorlesen“.
  \item Android-Automatisierung: Systemweite Szenen-Verknüpfungen (z. B. Xiaomi „Automatische Aufgaben“, Huawei „Szenen-Intelligenz“) sind am einfachsten. Tasker ist mächtig, hat aber eine steile Lernkurve – eher etwas für Technikbegeisterte.
  \item Grundidee: Erst die Handgriffe erkennen, die Sie täglich wiederholen, dann fragen: „Kann ein Auslöser mehrere Aktionen ersetzen?“ Beispiele und fortgeschrittene Varianten in Kapitel 23.
\end{itemize}

\booksection{Browser \& Suchmaschinen: Ihr Tor zum Internet}

Das Browser-Fenster ist Ihr Zugang zum Internet, die Suchmaschine der Kompass im Informationsmeer. Wer beide gut nutzt, spart sich viele Umwege.

\booksubsection{So arbeitet ein Browser:}

Nach Eingabe einer Adresse macht der Browser: (1) DNS-Abfrage nach der IP-Adresse zur Domain; (2) HTTP/HTTPS-Anfrage an den Server unter dieser IP; (3) der Server liefert den Seitenquelltext (HTML, CSS, JavaScript); (4) der Browser rendert daraus die sichtbare Seite.

\booksubsection{Empfohlene Browser:}

\begin{itemize}
  \item Microsoft Edge: Chromium-basiert, schnell und kompatibel. Mit integriertem Copilot-KI-Assistenten zum Zusammenfassen, Übersetzen und Chatten. Geräteübergreifende Synchronisierung – meine klare Empfehlung als Hauptbrowser.
  \item Safari: Standardbrowser auf Apple-Geräten, tief integriert, stromsparend und effizient. Guter Datenschutz, aber weniger Erweiterungen.
  \item Chrome: Weltweit am weitesten verbreitet, riesiges Erweiterungs-Ökosystem. Dafür speicherhungrig und am Handy eher stromfressend.
\end{itemize}

\booksubsection{Suchmaschine wählen:}

\begin{itemize}
  \item Bing: Microsofts Suchmaschine, tief mit Edge verknüpft. Die KI-Zusammenfassung liefert direkt Antworten statt nur Links – hohe Sucheffizienz. Bildersuche ebenfalls gut.
  \item Yandex: Russlands größte Suchmaschine, stark bei Bildersuche und Rückwärtssuche. Ideal, um ähnliche Bilder zu finden oder Objekte zu erkennen.
  \item Baidu: Die wichtigste Suchmaschine in China, aber mit viel Werbung – ein kritischer Blick ist nötig.
  \item Google: Weltweit die präziseste Suchmaschine; in China jedoch nur mit speziellem Netzwerkzugang nutzbar.
\end{itemize}

\booksubsection{Leitfaden für effiziente Recherche:}

\begin{itemize}
  \item KI-Suche nutzen: Bings KI-Zusammenfassung und Edge-Copilot beantworten Fragen direkt und destillieren die wichtigsten Fakten aus der Trefferflut – viel effizienter als Link für Link anklicken.
  \item Falschinformationen und Werbung erkennen: Treffer mit „Anzeige“ oder „Gesponsert“ sind bezahlte Platzierungen, keine organischen Ergebnisse. Quellen kritisch prüfen – offizielle Seiten und bekannte Medien zuerst.
  \item Apps korrekt herunterladen: In der Suche „App-Name + offizielle Website“ eingeben, z. B. „WeChat offizielle Website“, und Treffer mit offizieller Domain bevorzugen. Nicht von Drittanbieter-Downloadportalen installieren – dort wird oft Schadsoftware gebündelt.
  \item Offizielle Seiten erkennen: Windows 11 lädt man bei microsoft.com, Chrome bei google.com herunter. Ob die Domain zur Marke passt, ist die einfachste Prüfung.
\end{itemize}

\booksubsection{Domain-Wissen:}

Die Domain ist die „Hausnummer“ einer Website. Wer sie liest, erkennt Art und Herkunft der Seite:

\begin{itemize}
  \item Bedeutung häufiger Endungen: .com — kommerzielle Organisation, am häufigsten
  \item .cn — chinesische Länderdomain
  \item .org — gemeinnützige Organisation
  \item .net — Netzdienstleister
  \item .edu — Bildungseinrichtung
  \item .gov — Regierungsbehörde
  \item .xyz/.tech/.info — neuere generische Endungen, niedrige Hürden, Vorsicht geboten
\end{itemize}

Länder-/Regions-Domains: .jp — Japan

.kr — Südkorea

.uk — Vereinigtes Königreich

.de — Deutschland

.tw — Taiwan

.hk — Hongkong

Sicherheitshinweis: Seiten mit .gov und .edu sind in der Regel vertrauenswürdiger, weil die Vergabe streng ist. Billige Endungen wie .xyz oder .top werden häufig von Phishing-Seiten missbraucht – dort besondere Vorsicht.

\booksubsection{Werbeblocker:}

\begin{itemize}
  \item AdGuard Zielgruppe: Alle, die von Pop-ups, Banner- und In-App-Werbung genervt sind. Stärken: Werbeblockung auf Systemebene (nicht nur im Browser), filtert App-Banner und Tracking beim App-Start, spart Datenvolumen, beschleunigt Seiten; Premium für alle Geräte.
\end{itemize}

\booksubsection{Datenschutz-Browser:}

\begin{itemize}
  \item Brave Zielgruppe: Wer auf Privatsphäre achtet und nicht verfolgt werden will. Stärken: Chromium-basiert, blockiert ab Werk Werbung und Tracker, integrierter Tor-Modus, schneller als Chrome und speichersparender.
  \item DuckDuckGo Zielgruppe: Wer nicht will, dass eine Suchanfrage überall Werbung nach sich zieht. Stärken: Eigene Suchmaschine ohne Suchprotokoll, „Fire Button“ löscht alle Spuren mit einem Tap, einheitliche Nutzung über Plattformen hinweg.
\end{itemize}

\booksubsection{Browser im Überblick:}

\begin{center}
\begin{tabularx}{\linewidth}{|X|X|X|}
\hline
\textbf{Browser} & \textbf{Stärken} & \textbf{Zielgruppe} \tabularnewline
\hline
Microsoft Edge & Chromium-Kern + Copilot-KI: Seiten zusammenfassen, übersetzen, chatten & Produktive Büro-Nutzer, KI-Intensivnutzer (Hauptempfehlung) \tabularnewline
Safari & Tiefe Integration ins Apple-Ökosystem, stromsparend, starker Datenschutz & iPhone/Mac-Nutzer, die Akkulaufzeit und Systemkopplung schätzen \tabularnewline
Chrome & Weltweit meistgenutzt, größtes Erweiterungs-Ökosystem & Fortgeschrittene Nutzer mit vielen Add-ons, Fans geräteübergreifender Sync \tabularnewline
Brave & Blockiert ab Werk Werbung und Tracker, noch schneller & Datenschutzbewusste, die auch ohne Add-ons Werbung loswerden wollen \tabularnewline
DuckDuckGo & Eigene Suchmaschine ohne Verlauf, Spuren mit einem Tap löschen & Sensible Nutzer, die personalisierte Verfolgung ablehnen \tabularnewline
\hline
\end{tabularx}
\end{center}

\begin{tcbpractice}{Übung: Eine „Abendroutine“-Automatisierung anlegen}
\begin{enumerate}
  \item Auslöser erstellen: Bestimmte Uhrzeit (z. B. 22:30).
  \item Aktionen hinzufügen: Lautstärke senken > Nicht-stören einschalten > Entspannungs-Playlist abspielen > Helligkeit dimmen.
  \item Speichern und testen.
\end{enumerate}

\end{tcbpractice}

\booksection{Praktische Systemwerkzeuge: Aufräumen, Anti-Spam und Schutz}

Neben Produktivität und Gesundheit bestimmen drei Dinge den Alltag: „weniger gestört werden, genug Speicher, Betrüger bleiben draußen“. Zuerst die gängigen Helfer, dann die integrierten Blockierfunktionen des Smartphones – viele haben sie seit Jahren nie eingeschaltet.

\booksubsection{System bereinigen und Speicher optimieren}

\begin{itemize}
  \item Qingzhuo – Für wen geeignet: Android-Nutzer, die Cache, Datei-Reste und leere Ordner unkompliziert aufräumen wollen, ohne sich mit einem aufgeblähten „Telefon-Manager“-Paket herumzuschlagen. Das Besondere: klein und werbefrei, mit Tiefenscan für App-Reste, Auffinden großer Dateien und Analyse der Speicherbelegung; zurückhaltender als kommerzielle „Beschleunigen auf Knopfdruck“-Apps – ideal für technisch versierte Nutzer zur regelmäßigen Pflege. Achtung: Prüfe vor dem Aufräumen die Pfade; bei WeChat, Fotoalbum und Dokumentenordnern nicht blind alles auswählen. Für Cache auf Systemebene zuerst die Speicherreinigung des Herstellers nutzen – Drittanbieter-Tools nur als Ergänzung.
  \item Integrierte Speicherbereinigung Zielgruppe: Alle, besonders ältere Nutzer und Technikmuffel. Stärken: Android [Einstellungen] > [Speicher] > [Aufräumen]; iOS [Einstellungen] > [Allgemein] > [iPhone-Speicher]. Das System kennt die eigene App-Struktur am besten, geringstes Risiko, Wichtiges zu löschen – daher erste Wahl.
  \item WeChat/QQ-Speicherverwaltung Zielgruppe: Vielchatter. Stärken: WeChat [Ich] > [Einstellungen] > [Allgemein] > [Speicherplatz], fein nach Chat und Dateityp aufräumen. Sicherer als externe „WeChat-Cleaner“.
  \item Prinzip: Was System oder App selbst kann, nicht an Dritte geben. Was nach Pfad/Typ wählbar ist, nicht „alles auf einen Schlag“ löschen. Ziel ist freier Speicher – nicht das versehentliche Mitlöschen wichtiger Daten.
\end{itemize}

\booksubsection{Anti-Spam und Sicherheits-Apps}

\begin{itemize}
  \item National Anti-Fraud Center Zielgruppe: Alle, besonders ältere Menschen, Finanzberufe, Vielbesteller und Anrufunbekannter. Stärken: Offiziell, mit Anrufwarnung, SMS-Risikohinweis und verdächtige-URL-Prüfung. Angebunden an die Anti-Betrug-Datenbank der Polizei – das offiziellste Glied der „drei Verteidigungslinien“.
  \item Sicherheitscenter / Handy-Manager der Marke Zielgruppe: Nutzer der jeweiligen Marke. Stärken: Huawei Mobile Manager, Xiaomi Security Center, OPPO/vivo Mobile Manager usw. Integrieren Spam-Schutz, Virensuche, Rechteverwaltung und Zahlungsschutz. Meist kostenlos und systemnah – keine zweite kommerzielle Sicherheits-App installieren, sonst drohen Rechtekonflikte.
  \item Anbieter-Anti-Spam Zielgruppe: Wer von Massenanrufen und KI-Sprachanrufen bombardiert wird. Stärken: China Mobile: „KTFSR“ an 10086 senden; China Telecom: „Tianyi Anti-Spam“ folgen; China Unicom: im „Unicom Mobile Manager“ einstellen. Die Blockierung läuft netzseitig – selbst wenn das Handy nicht klingelt, sind viele Spam-Nummern draußen.
  \item AdGuard Zielgruppe: Alle, die von Pop-ups, Banner- und In-App-Werbung genervt sind. Stärken: Systemweite Werbeblockung (nicht nur Browser), filtert Banner und Tracking beim App-Start, spart Datenvolumen, beschleunigt Seiten; Premium für alle Geräte.
  \item Vorsicht bei „Cleaner/Anti-Spam“-Paketen: Manche Drittanbieter-Manager laufen dauerhaft im Hintergrund und verlangen wiederholt Kontakt- und SMS-Rechte. Ihre Blockierkraft übertrifft das System oft nicht, kosten aber Akku und Privatsphäre. Zuerst die integrierten Funktionen nutzen, dann offizielle Anti-Betrug-Apps.
\end{itemize}

\textbf{Auswahlprinzipien für Werkzeug-Apps}\par

\begin{itemize}
  \item System zuerst: Blockieren, Aufräumen, Rechte – zuerst die Schalter im Smartphone selbst.
  \item Offizielle Ergänzung: National Anti-Fraud Center und Anbieter-Spamschutz sind offizielle Kanäle und dürfen sich stapeln.
  \item Drittanbieter sparsam: Pro Funktion eine App – vermeidet Rechtekämpfe und Hintergrundkonkurrenz.
  \item Weniger ist mehr: Je mehr Werkzeuge, desto größer die Angriffsfläche für Benachrichtigungen und Rechte – und desto eher stört „das Werkzeug selbst“.
\end{itemize}

\booksubsection{Integrierte Spam-Sperren: So aktivieren Sie sie je Marke}

Die Einstiegspfade je Marke folgen unten. Namen können je Systemversion leicht abweichen; in den Einstellungen einfach nach „Spam“, „Blockieren“ oder „Unbekannte Anrufe“ suchen. Ergänzend zu Kapitel 8 „Belästigungsanrufe“: Hier geht es um die Schalter, dort um Nummernkennzeichnung und Einspruch bei Fehlmarkierungen.

\begin{center}
\begin{tabularx}{\linewidth}{|X|X|X|}
\hline
\textbf{Plattform} & \textbf{Einstieg} & \textbf{Empfohlene Optionen} \tabularnewline
\hline
iPhone / iOS & {}[Einstellungen] > [Telefon] / [Nachrichten] & Unbekannte Anrufe stumm schalten; blockierte Kontakte; in Nachrichten „Unbekannte Absender filtern“ \tabularnewline
Xiaomi / Redmi & Telefon-App > Einstellungen > Spam-Schutz; oder Security Center & Intelligente Blockierung, „einmal klingeln und auflegen“ bei Unbekannten, SMS-Keyword-Filter \tabularnewline
Huawei / Honor & Telefon > Spam-Schutz; Mobile Manager > Spam-Schutz & Anruferkennung, Schwarze Liste, intelligente SMS-Filterung \tabularnewline
OPPO / realme / OnePlus & Telefon-App > Spam-Schutz; oder Gerät > Spam-Schutz & Anrufblockierung, Betrugserkennung, SMS-Filterregeln \tabularnewline
vivo / iQOO & Telefon > Spam-Anrufe blockieren; i-Manager & Spam-Schutz, Betrugswarnung, Nachrichtenfilter \tabularnewline
Samsung & Telefon > Einstellungen > Nummern blockieren / Anrufer-ID \& Spam & Unbekannte Nummern blockieren, Spam-Erkennung (je Region unterschiedlich) \tabularnewline
\hline
\end{tabularx}
\end{center}

\begin{tcbinfo}{Drei Ebenen kombinieren}
Empfehlung: „System-Blocker + Anbieter-Anti-Spam + National Anti-Fraud Center“. Das System markiert und blockiert, der Anbieter fängt Massenanrufe netzseitig ab, das Anti-Fraud Center liefert offizielle Warnungen. Alle Ebenen nicht auf Maximum stellen – sonst werden auch Lieferdienste und Arztanrufe blockiert. Kontakte und „Zuletzt kontaktiert“ freigeben.

\end{tcbinfo}

\begin{tcbpractice}{Übung: Spam-Schutz heute noch einschalten}
\begin{enumerate}
  \item Telefon-App-Einstellungen öffnen, „Spam-Schutz / Unbekannte Anrufe stumm“ finden, intelligente Blockierung aktivieren.
  \item Dasselbe bei den Eltern erledigen: gleicher Weg + Anbieter-Anti-Spam (bei China Mobile „KTFSR“ an 10086).
  \item „National Anti-Fraud Center“ aus dem App Store installieren, registrieren und Anrufwarnung einschalten.
  \item Testnummer anrufen (oder Blockierprotokoll prüfen), um die Funktion zu bestätigen.
\end{enumerate}

\end{tcbpractice}

\booksection{Wetterwarnungen, Naturgefahren und Erdbebenwarnung}

Spam-Schutz stoppt Telefonbetrug, Warnsysteme stoppen Naturgefahren und Extremwetter. Viele installieren nur eine Wetter-App, lassen Benachrichtigungen und Erdbebenwarnung aus – dabei können die Sekunden vor Starkregen, Taifun oder Beben Leben retten.

\booksubsection{Wetterwarnungen: Starkregen, Taifun, Hitze}

\begin{itemize}
  \item Offizielle Kanäle zuerst: Die Warnungen von China Meteorological Administration, dem Zentralen Wetterdienst und den Provinzdiensten sind am verlässlichsten. „China Weather“ auf WeChat/Weibo folgen oder offizielle Apps wie „China Weather“ installieren.
  \item System-Wetter-Apps: Huawei, Xiaomi, OPPO, vivo, Apple „Wetter“ usw. pushen meist lokale Wetterdienst-Warnungen. In den App-Einstellungen prüfen, ob „Wetterwarnungen / Unwetter“ aktiviert sind.
  \item Minutengenaue Regenprognosen wie Caiyun: Gut für Pendeln oder den Schulweg – „Brauche ich in den nächsten 30 Minuten einen Regenschirm?“ Für offizielle Warnungen gilt weiterhin der Wetterdienst.
  \item Warnstufen (allgemein): Blau: Aufmerksam sein, Reise wie geplant, Updates beachten.
  \item Gelb: Erhöhte Vorsicht, unnötige Outdoor-Aktivitäten reduzieren.
  \item Orange: Ernste Lage, möglichst drinnen bleiben, Gefahrenzonen meiden.
  \item Rot: Besonders ernst. Anweisungen zu Schließungen und Evakuierung folgen; keine überfluteten Straßen, Wildbäche oder Sturmpassagen riskieren.
\end{itemize}

Häufige Warnarten: Starkregen, Schneesturm, Taifun, Hitze, Kältewelle, Sturm, Blitz, Hagel, Glatteis, Nebel, Dürre, geologische Risiken.

\booksubsection{Naturgefahren: Mehr als nur „Regen“}

\begin{itemize}
  \item Wer warnt: Wetterdienste für Wetterlagen; geologische Gefahren und Sturzfluten oft gemeinsam mit Umwelt-, Wasser- und Katastrophenschutzbehörden. Kommunen nutzen Notfallfunk, SMS und Kanäle wie 12379.
  \item Warnungen mehrkanalig empfangen: Systembenachrichtigungen + Wetter-App.
  \item WeChat: Lokalen „Katastrophenschutz“ und „Wetterdienst“ folgen.
  \item 12379 – nationale Warnkanäle (SMS/Dienst, je nach Region freigeschaltet).
  \item Fernsehen, Notfallfunk, Dorf-/Viertellautsprecher – in Bergregionen und zur Hochwasserzeit nicht nur aufs Handy schauen.
\end{itemize}

Draußen unterwegs: Bei Taifun nicht an Reklametafeln oder unter Bäumen gehen; bei Starkregen keine Tiefgaragen, Tunnel oder Unterführungen; in den Bergen bei Warnung vor geologischen Risiken sofort Hänge und Schluchten verlassen.

Verzahnung mit dem Notfallhandbuch: Bei Kontaktabbruch, leerem Akku oder Standortnotruf siehe Kapitel 32 (Erdbeben/Katastrophen) und Kapitel 7 (Satellitenkommunikation ohne Netz).

\booksubsection{Erdbebenwarnung (EEW) ist nicht „Erdbebenprognose“}

Die Erdbebenwarnung (EEW) nutzt, dass elektromagnetische Signale schneller sind als Bebenwellen und Stationen die schnellen P-Wellen früher erfassen. So erhalten betroffene Regionen Sekunden bis Dutzende Sekunden Vorlauf, bevor die zerstörerischen S-Wellen eintreffen. Sie sagt nicht voraus, ob ein Beben kommt, schafft aber Zeit zum Schutz, zum Abschalten von Feuerquellen oder zum Verlassen des Aufzugs. Je näher am Epizentrum, desto kürzer die Vorwarnzeit – im Extremfall bleibt kaum Zeit.

\begin{itemize}
  \item Offizielle und branchenweite Lösungen: In Festlandchina gehören dazu Fähigkeiten des China Earthquake Networks sowie Warnsysteme von Instituten wie dem Chengdu Institute of Disaster Prevention. Manche Geräte bringen Warnfunktionen vorinstalliert oder empfangen sie online.
  \item Drittanbieter-Apps: Im App Store nach „Erdbebenwarnung“ und ähnlichen offiziellen/Instituts-Apps suchen. Vor der Installation den Entwickler prüfen – vor „Fake-Warn“-Apps, die unnötige Rechte verlangen.
  \item Nach einer Warnung: Drinnen in der Nähe Schutz suchen (unter Tischen, an tragenden Wänden, an stabilen Möbeln), Kopf und Nacken schützen; von Fenstern, Außenwänden und hohen Möbeln weg; nicht blind in Treppenhäuser oder aus dem Fenster; im Aufzug auf der nächsten Etage aussteigen; draußen offene Flächen aufsuchen, fern von Leitungen, Reklame und maroden Gebäuden. Nach dem Beben Gas und Strom prüfen, Nachbeben beachten.
  \item Nicht stumm lassen: Die Erdbebenwarnung muss den Stumm-/Nicht-stören-Modus durchbrechen dürfen. Nachts das Handy griffbereit legen, Klingellautstärke nicht auf 0.
\end{itemize}

\booksubsection{Integrierte Erdbebenwarnung: je Marke aktivieren}

Die Pfade können je Systemversion leicht abweichen; in den Einstellungen nach „Erdbeben“, „Warnung“ oder „Sicherheitswarnung“ suchen. iOS hat in China nur begrenzte integrierte Unterstützung – meist ist eine verlässliche Drittanbieter-Warn-App nötig, mit Standort- und Benachrichtigungsrechten.

\begin{center}
\begin{tabularx}{\linewidth}{|X|X|X|}
\hline
\textbf{Plattform} & \textbf{Häufiger Einstieg} & \textbf{Hinweis} \tabularnewline
\hline
Xiaomi / Redmi & Mobile Manager / Security Center > Erdbebenwarnung; oder in Einstellungen nach „Erdbeben“ suchen & Manche Systemversionen bringen einen Warnservice mit; Internet und Standort nötig \tabularnewline
Huawei / Honor & Einstellungen > Sicherheit > Notfallwarnungen / Erdbebenwarnung & Kann „Notfallwarnungen“ heißen, enthält u. a. Erdbeben \tabularnewline
OPPO / realme / OnePlus & Einstellungen > Sicherheit (SOS-Notruf) > Erdbebenwarnung & Je Modell und Region unterschiedlich verfügbar \tabularnewline
vivo / iQOO & Wetter-App oder Einstellungen > Sicherheit/Mehr > Erdbebenwarnung & Maßgeblich ist das Menü der installierten Systemversion \tabularnewline
iPhone / iOS & Oft kein vollständiger integrierter Schalter; verlässliche Erdbeben-Warn-App installieren und Benachrichtigungen erlauben & App braucht Hintergrundbetrieb und Standort; nicht-stören darf ihre Warnung nicht unterdrücken \tabularnewline
Samsung u. a. & Einstellungen > Sicherheit \& Notfälle / Drahtlos-Notfallwarnungen & Abdeckung von „Notfallwarnungen“ variiert je Region und Anbieter \tabularnewline
\hline
\end{tabularx}
\end{center}

\begin{tcbwarn}{Häufige Missverständnisse zur Erdbebenwarnung}
\begin{itemize}
  \item Warnung $\neq$ Prognose: Sie sagt nicht, „ob morgen ein Beben kommt“, sondern meldet nur ein bereits geschehenes Beben vor Ankunft der Wellen.
  \item Nicht überall viele Sekunden Vorlauf: Nahe am Epizentrum bleiben oft nur Sekunden oder gar keine Zeit; vor allem die Umgebung profitiert.
  \item Benachrichtigungen aus = App sinnlos: Selbst die beste Warnung kommt nicht an, wenn sie stumm geschaltet, im Hintergrund gekillt oder ohne Standort läuft.
  \item Fälschliche Apps: Vor „Erdbebenprognose“-Software, die Kontakte und SMS-Rechte verlangt, warnen. Integrierte Lösungen oder seriöse Anbieter bevorzugen.
\end{itemize}

\end{tcbwarn}

\begin{tcbinfo}{Warnungen sinnvoll kombinieren}
Alltagsempfehlung: „System-Wetterwarnung + offizielle Wetter-/Notfallkanäle + Erdbebenwarnung (System oder seriöse App)“. Wie beim Spam-Schutz gilt: nicht dutzende gleiche Apps installieren. Wichtig sind Benachrichtigungsrechte, Standorterlaubnis und die Erlaubnis, den Nicht-stören-Modus zu durchbrechen. Nach dem Spam-Schutz für die Eltern auch Wetter- und Erdbebenwarnung prüfen.

\end{tcbinfo}

\begin{tcbpractice}{Übung: Warn-Einstellungen in 10 Minuten prüfen}
\begin{enumerate}
  \item System-Wetter-App öffnen, in die Einstellungen und „Wetterwarnungen / Unwetter“ aktivieren.
  \item In den Systemeinstellungen nach „Erdbeben“ oder „Warnung“ suchen und die Warnung einschalten (Android). Auf dem iPhone eine seriöse Erdbeben-Warn-App installieren und Benachrichtigungen erlauben.
  \item Dem lokalen Wetterdienst oder Katastrophenschutz in WeChat folgen und Push erlauben.
  \item Prüfen, dass die Warntöne vom Nicht-stören-Modus nicht komplett unterdrückt werden; dasselbe am Handy der Eltern.
\end{enumerate}

\end{tcbpractice}

\booksection{Geräteübergreifend: PC-Fernsteuerung, Schnelltransfer und Kontosicherheit}

Ein Smartphone läuft besonders gut, wenn es mit PC und Tablet ein Team bildet. Hier ergänzend die wichtigsten Helfer: Fernsteuerung des PCs, KI-Büro und Remote-Zusammenarbeit, Live-Wallpapers, LAN-Transfer, geräteübergreifende Zwischenablage, TOTP-2FA sowie Papier-Scans und Terminplanung. Systematischer zu Mehrgeräte-Nutzung: Kapitel 17.

\booksubsection{Smartphone steuert den PC: ToDesk / Sunlogin}

\begin{itemize}
  \item Sunflower Remote Control – Für wen geeignet: alle, die von unterwegs auf ihren Windows- oder Mac-Rechner zu Hause oder im Büro zugreifen wollen; auch beliebt, um Älteren aus der Ferne beim Handy zu helfen (siehe Kapitel 27). Das Besondere: ausgereiftes Angebot aus China – per Handy-App den PC-Desktop steuern, Dateien übertragen, den Rechner ein- und ausschalten (teilweise mit Zusatzhardware); die Gratisversion reicht für den leichten Privatgebrauch meist aus.
  \item ToDesk Zielgruppe: Stabile, latenzarme Remote-Arbeit und -Wartung. Stärken: Verbindung über Gerätecode + Passwort, Handy steuert PC, PC steuert Handy (je Version), Dateitransfer und Mehrbildschirm gut gelöst; in China oft gute Knotenleistung.
  \item Grundablauf: Client am PC installieren und online halten > App am Handy > Sitzung über Geräte-ID + Verbindungspasswort (oder gleiches Konto) > Desktop per Touchpad/Gesten am Handy steuern.
  \item Sicherheitsregeln (unbedingt beachten): Wenn Fremde sich als „Support/Behörde/Anbieter“ ausgeben und Sie auffordern, Sunlogin oder ToDesk zu installieren, um Ihr Gerät fernzusteuern: immer ablehnen – das ist ein typischer Betrugsversuch (siehe Kapitel 27).
  \item Eigene Nutzung: Zwei-Faktor einschalten, Standardpasswort ändern, bei Nichtnutzung abmelden oder unbeaufsichtigten Zugriff abschalten. Firmenrechner nach internen Sicherheitsregeln nutzen.
  \item Im öffentlichen Raum darauf achten, dass niemand Bildschirm und Passwörter mitliest.
\end{itemize}

\booksubsection{KI-Büro: Trae / WorkBuddy / Dumate}

\begin{itemize}
  \item Trae Zielgruppe: Programmieren und Dokumenten-Zusammenarbeit, freundlich für chinesische Umgebungen. Stärken: KI-IDE/-Assistent von ByteDance, SOLO-Modus erzeugt Projekte aus einem Satz. Effizient beim Coden und Texten. Läuft der PC, lassen sich Aufgaben am Mobilgerät per Fernzugriff oder offizieller Begleit-App verfolgen (Produktstand prüfen). Details in Kapitel 41 „KI-Programmierwerkzeuge im Überblick“.
  \item WorkBuddy Zielgruppe: Berufstätige, die KI bei Meetingprotokollen, Mails, Berichten und Office-Dokumenten unterstützen wollen. Stärken: KI-Assistent fürs Büro, gut für „Entwerfen – Glätten – Zusammenfassen“. Manche Workflows starten am Handy, schwere Arbeit läuft am PC.
  \item Dumate Zielgruppe: Teams und Einzelpersonen, die Kollaboration, Dokumente und Wissensmanagement vereinen wollen. Stärken: KI-gestütztes Office mit Fokus auf Materialbündelung und gemeinsames Bearbeiten. Am Handy sichten und leicht bearbeiten; komplexe Layouts am großen Bildschirm.
  \item „Handy steuert PC“ verstehen: Viele KI-Bürofunktionen laufen am PC oder im Browser. Mit Sunlogin/ToDesk zurück auf den Rechner greifen – oder die eigene App/des Anbieters Webseite nutzen – und unterwegs KI-Aufgaben „remote“ steuern. Faustregel: Rechenlast und lange Dokumente an den PC, Fortschritt und kleine Änderungen ans Handy.
  \item Datenschutz: Vor dem Upload von Verträgen, Ausweisen oder Kundendaten die Datenschutzbestimmungen und Compliance-Anforderungen prüfen. Firmdaten bevorzugt mit firmenintern freigegebenen Tools bearbeiten.
\end{itemize}

\booksubsection{Live-Wallpaper: miHoYo Desktop (Handy + PC)}

\begin{itemize}
  \item Rengong Desktop – Für wen geeignet: Fans von miHoYo-Charakteren, die kostenlose animierte Hintergrundbilder wollen. Das Besondere: eine kostenlose App für animierte Hintergründe von miHoYo, für Handy und PC; die Live2D-/Animationsdarstellungen der Charaktere dienen als Desktop-Bild, auf manchen Geräten/Systemen auch als Sperrbildschirm- oder Desktop-Hintergrund.
  \item Strom und Performance: Live-Wallpapers verbrauchen mehr Akku und GPU als statische Bilder. Bei wenig Strom oder älteren Geräten nur während des Ladens nutzen oder die Bildrate senken (falls die App es anbietet).
  \item Installationsquelle: Offizielle Website oder seriöser App Store – keine „gecrackten“ Pakete mit Trojanern. Bei der PC-Version Installationspfad und Autostart prüfen.
\end{itemize}

\booksubsection{LAN-Schnelltransfer: LocalSend / IPMsg / AirDrop}

\begin{itemize}
  \item LocalSend Zielgruppe: Gemischte Android/Windows/iOS/Mac/Linux-Setups, plattformübergreifend ohne Cloud. Stärken: Open Source, Direktverbindung im LAN, Fotos/Videos/Dokumente im selben WLAN austauschen, ohne Kabel und ohne Herstellerserver – Tempo hängt vom Router ab.
  \item Feige Chuanshu – Für wen geeignet: Windows-PCs im Büro, klassische Arbeitsumgebungen im lokalen Netzwerk. Das Besondere: ein etabliertes LAN-Tool für Übertragung und Chat, ideal für große Dateien im internen Netz; nur von offiziellen Quellen herunterladen und im Firmennetz die IT-Richtlinien beachten.
  \item AirDrop Zielgruppe: iPhone-, iPad- und Mac-Nutzer. Stärken: Teilen im Apple-Ökosystem per Tap, Bluetooth-Erkennung + WLAN-Direktverbindung, ohne Datenvolumen. „Nur Kontakte“ senkt Belästigung durch Fremde.
  \item Ergänzend: Mi Share, Huawei Share, vivo Share usw. sind innerhalb derselben Marke sehr schnell. WeChat „Dateiübertragungs-Assistent“ ist bequem, aber dateigrößen- und gerätebeschränkt; für große Dateien eher LocalSend/AirDrop.
  \item Nutzungstipp: Beide Geräte im selben WLAN. Auf öffentlichem WLAN sensible Dateien besser über den privaten Hotspot senden. Nach dem Transfer beim Empfänger die Vollständigkeit prüfen.
\end{itemize}

\booksubsection{Zwischenablage geräteübergreifend: KDE Connect / Handy-PC-Transfer}

\begin{itemize}
  \item Universelle Zwischenablage (Apple) Zielgruppe: iPhone + Mac. Stärken: Gleiches Apple-ID und Handoff: Am Handy kopieren am Mac einfügen (und umgekehrt). Text und manche Dateien funktionieren.
  \item KDE Connect Zielgruppe: Android + Windows/Linux/Mac, Open-Source-Zwischenablage, Spiegelung von Benachrichtigungen, Fernsteuerung. Stärken: Nach LAN-Kopplung Zwischenablage synchronisieren, Handy-Benachrichtigungen an den PC, Handy als Präsentations-Fernbedienung/Touchpad. Daten gehen nicht über Drittanbieterkonten.
  \item Herstellereigene „Handy-PC-Übertragung“: Huawei Multi-Screen Collaboration/Super Device, Xiaomi Interconnect, vivo Share, OPPO Connect, Honor Connect usw. Oft mit plattformübergreifender Zwischenablage und Drag-and-drop – innerhalb einer Marke am besten.
  \item Windows + Android, offizieller Weg: „Phone Link / Link to Windows“ (Nachfolger von Your Phone) zeigt Handybenachrichtigungen und Fotos am PC und unterstützt Kopieren/Einfügen in unterstützten Szenarien. Bei Samsung und anderen tiefer integriert.
  \item Sicherheitshinweis: In der Zwischenablage können Passwörter und Codes zurückbleiben. Auf geteilten Rechnern globale Sync vorsichtig nutzen, nach dem Gebrauch keine erneuten Einfügevorgänge in Passwortfeldern. Manche Tastatur-Apps lesen die Zwischenablage – Rechte beachten.
\end{itemize}

\booksubsection{TOTP-2FA: Aegis / Microsoft Authenticator}

\begin{itemize}
  \item Warum ein Authenticator statt SMS: SMS-Codes können über SIM-Tausch gekapert werden. TOTP-Codes entstehen lokal zeitbasiert, ohne Handynetz, und sind resistenter gegen Phishing und Abfangen (Trotzdem vor Malware und Screenshots hüten).
  \item Aegis Authenticator Zielgruppe: Android-Nutzer, die Open Source und lokale Kontrolle schätzen. Stärken: Kostenlos und quelloffen, verschlüsselte Backups, Import/Export, aufgeräumte Oberfläche, Token-Liste per Fingerabdruck/Gesicht freischaltbar.
  \item Microsoft Authenticator Zielgruppe: Microsoft-/Office-365-/Firmenkonten, unter iOS und Android. Stärken: Push-Bestätigung für Anmeldungen, Passwort-Überwachung und Konten-Backup (je Produktversion), enge Verbindung mit Windows und Unternehmens-Identitäten.
  \item Google Authenticator / 2FAS usw.: Auch im jeweiligen Ökosystem verbreitet. Wichtig ist ein verschlüsseltes Backup, damit Tokens beim Geräte-Upgrade nicht verloren gehen.
  \item Beim Geräte-Upgrade: Daten vom alten Gerät migrieren oder „Ersatzcodes/Wiederherstellungsschlüssel“ aus dem Dienst nutzen. Ersatzcodes ausdrucken und an zwei Orten verwahren – nicht nur auf demselben Handy.
  \item Mit Passwortmanagern: Manche Tresore speichern auch TOTP – praktisch, aber das Risiko bündelt sich. Für hochwertige Konten empfiehlt sich die Trennung „Passwortmanager + eigener Authenticator“. Details in Kapitel 14.
\end{itemize}

\booksubsection{Empfohlene Kombination: Ein Tag im Überblick}

\textbf{Beispiel-Workflow}\par

\begin{itemize}
  \item Morgens: TickTick öffnen, Tagesplan checken > Wetter/Warnungen überfliegen > am PC Trae/Dokumente starten, Handy beiseitelegen.
  \item Tagsüber: Papierunterlagen mit einer Scan-App in PDF umwandeln, per LocalSend/AirDrop an den PC senden; lange Texte am PC mit KI-Bürotools entwerfen.
  \item Unterwegs: Mit ToDesk/Sunlogin auf den PC zugreifen und weiterarbeiten; wichtige Logins über Aegis/Microsoft Authenticator absichern.
  \item Abends: In TickTick abhaken; bei Bedarf mit miHoYo Desktop das Wallpaper wechseln und durchatmen.
\end{itemize}

\begin{tcbwarn}{Faustregeln für Fernzugriff und Kontosicherheit}
\begin{itemize}
  \item Jemand verlangt, Fernsteuerungssoftware zu installieren > auflegen und selbst den offiziellen Support anrufen.
  \item Bestätigungs- und Einmalpasswörter niemandem mitteilen – auch nicht „Paketdienst“ oder „Support“.
  \item Vor Fernzugriff auf Firmenrechnern oder Installation von KI-Software zuerst die IT fragen.
  \item Dateiübertragung und Fernanmeldung in öffentlichen Netzen besonders kritisch behandeln.
\end{itemize}

\end{tcbwarn}

\begin{tcbpractice}{Übung: Die geräteübergreifende Kette einmal durchspielen}
\begin{enumerate}
  \item Handy und PC ins selbe WLAN, per LocalSend (oder AirDrop) ein Foto austauschen.
  \item In TickTick die Aufgabe „Heute 20 Uhr: Scans sortieren“ anlegen und prüfen, ob sie auf Handy und PC erscheint.
  \item Für E-Mail oder WeChat 2FA aktivieren (falls unterstützt), mit Aegis oder Microsoft Authenticator verknüpfen und Ersatzcodes sichern.
  \item Sunlogin oder ToDesk am eigenen PC installieren, einmal aus dem eigenen Handy verbinden, dann abmelden oder Passwort härten – Ablauf üben, aber keine Sicherheitslücke offen lassen.
\end{enumerate}

\end{tcbpractice}
